\clearpage\section*{References}
{\footnotesize\raggedright
\begin{itemize}[leftmargin=1.1em,itemsep=0.3em,parsep=0pt]
\item \textit{Missale Romanum}, editio typica (Vatican City: Typis Polyglottis Vaticanis, 1962), \latin{Dominica decima sexta post Pentecosten}, pp.~397--398, nos.~1592--1601. \sourceurl{https://media.churchmusicassociation.org/pdf/missale62.pdf}{CMAA facsimile}. \textit{Missale Romanum} (Ratisbon: Pustet, 1862), same Sunday, pp.~339--341, \sourceurl{https://archive.org/download/bub_gb_E7sPAAAAIAAJ/bub_gb_E7sPAAAAIAAJ.pdf}{facsimile}.
\item \textit{The Roman Missal translated into the English language for the use of the laity} (Philadelphia: Eugene Cummiskey, 1861), XVI.~Sunday after Pentecost, Collect, Secret and Postcommunion; \sourceurl{https://archive.org/download/romanmissaltran00churgoog/romanmissaltran00churgoog.pdf}{facsimile}. Unnamed historical translator; conclusions remain abbreviated.
\item English College at Douay and Rheims; Richard Challoner, reviser, \textit{The Bible, Douay-Rheims, Complete}, \sourceurl{https://www.gutenberg.org/ebooks/1581}{Gutenberg eBook 1581}, electronic text released 1 December 1998, updated 23 September 2023; electronic production by Dennis McCarthy and Tad Book. Pss.~39,69,70,77,85,97,101; Eph.~1--6; Luke~13--15, 1:46--55, 2:22--39, 18:9--14, 22:24--30; Exod.~23,34; Deut.~22; Prov.~25; Isa.~52,60; John~14; Col.~1--2; Heb.~1,10. The eBook number identifies no single printed exemplar.
\item \textit{Biblia Sacra Vulgatae Editionis Sixti V et Clementis VIII}, Clementine text, eBible \sourceurl{https://ebible.org/latVUC/}{latVUC edition}: appointed Psalm loci, Eph.~3 and Luke~14, for comparison with the Missal's adaptations.
\item Augustine, \textit{Enarrationes in Psalmos}, in \textit{Nicene and Post-Nicene Fathers}, first series VIII, English, \sourceurl{https://ccel.org/ccel/s/schaff/npnf108/cache/npnf108.txt}{CCEL complete offered digital volume}: Latin Ps.~39, NPNF ``Psalm XL,'' §§21--25; Ps.~70, ``Psalm LXXI,'' §§18--21; Ps.~85, ``Psalm LXXXVI,'' §§1--3,5,7; Ps.~97, ``Psalm XCVIII,'' §1; Ps.~101, ``Psalm CII,'' §§16--18. The edition contains explicit abridgments.
\item Augustine, \textit{Epistula}~140 to Honoratus / \textit{De gratia Novi Testamenti}, 25.62--26.64, with §§60--67; \sourceurl{https://www.augustinus.it/latino/lettere/lettera_141_testo.htm}{Augustinus.it Latin digital edition}. The page identifies Letter~140 despite the delivery filename.
\item John Chrysostom, \textit{Homilies on Ephesians} VII, on Eph.~3:13--21 with opening at 3:8, \textit{NPNF}, first series XIII, \sourceurl{https://ccel.org/ccel/s/schaff/npnf113/cache/npnf113.txt}{CCEL English edition}, physical lines~8030--8176.
\item Theodoret of Cyrus, \textit{Interpretatio in Psalmos}, PG~80 (Paris, 1860), paired Greek and Latin: Ps.~39:13--17, cols.~1159A--1160D; Ps.~70:15--19, 1423B--1426C; Ps.~85:1--5, 1553C--1556B; Ps.~97:1, 1657C--1658D; Ps.~101:14--19, 1679B--1682B. \sourceurl{https://upload.wikimedia.org/wikipedia/commons/f/fd/Patrologia_Graeca_Vol._080.pdf}{Facsimile}.
\item Cyril of Alexandria, \textit{Commentary on Luke}, trans. R.~Payne Smith from Syriac (Oxford, 1859), Sermons CI--CII, pp.~471--479. \sourceurl{https://www.tertullian.org/fathers/cyril_on_luke_10_sermons_99_109.htm}{Roger Pearse's web edition}, sermon division 99--109. The son/ox text is distinct from the Missal's ass/ox.
\item Ambrose, \textit{Expositio evangelii secundum Lucam} VII.195, PL~15, cols.~1752A--B, \sourceurl{https://la.wikisource.org/wiki/Expositio_evangelii_secundum_Lucam/VII}{Latin Wikisource text}, with preceding and following context.
\item Thomas Aquinas, \textit{Super Epistolam B.~Pauli ad Ephesios lectura}, chapter~3, lectures~4--5, units~87804--87805, \sourceurl{https://www.corpusthomisticum.org/cep.html\#87804}{Corpus Thomisticum Latin}, Turin 1953 base, Roberto Busa transcription revised by Enrique Alarcón. \textit{Super Psalmo}~39 nn.~6--7, unit~87202, reportatio of Reginald of Piperno, Parma 1863 base, \sourceurl{https://www.corpusthomisticum.org/cps31.html\#87202}{Corpus Thomisticum Latin}.
\item Robert Bellarmine, \textit{A Commentary on the Book of Psalms}, trans. John O'Sullivan (Dublin: James Duffy, 1866): Ps.~LXX:16--18, pp.~212--213; LXXXV:1,3--5, pp.~259--260; XCVII:1, p.~306; CI, Gentiles and built-Sion lemmas, p.~317 (edition's comments 15--16). \sourceurl{https://upload.wikimedia.org/wikipedia/commons/9/9d/Commentaryonbook0000bell.pdf}{Facsimile}.
\item Francis de Sales, \textit{Introduction à la vie dévote} III.5, ``De l'humilité plus intérieure,'' especially the paragraph beginning \latin{Nous disons souvent que nous ne sommes rien} (Lyon/Paris: Perisse Frères, 1832), \sourceurl{https://www.gutenberg.org/ebooks/53540.txt.utf-8}{Gutenberg eBook 53540}, French.
\item Benedict XVI, \sourceurl{https://www.vatican.va/content/benedict-xvi/en/angelus/2010/documents/hf_ben-xvi_ang_20100829.html}{Angelus, 29 August 2010}, first two Gospel paragraphs, overlapping postconciliar Luke~14:1,7--14; and \sourceurl{https://www.vatican.va/content/benedict-xvi/en/audiences/2009/documents/hf_ben-xvi_aud_20090114.html}{General Audience, 14 January 2009}, ``Saint Paul (18): The theological vision of the Letters to the Colossians and Ephesians,'' mystery paragraph. Official English, paraphrased.
\item H.~A. Wilson, ed., \textit{The Gelasian Sacramentary} (Oxford: Clarendon Press, 1894), III.12, p.~231, \sourceurl{https://archive.org/download/gelasiansacrame00wilsgoog/gelasiansacrame00wilsgoog.pdf}{facsimile}; \textit{The Gregorian Sacramentary under Charles the Great}, HBS~49 (London: Harrison and Sons, 1915), pp.~133--135,173 note~c,174, introduction pp.~xvii--xviii,xxx,xli,xliv--xlv, \sourceurl{https://archive.org/download/gregoriansacrame00cath/gregoriansacrame00cath.pdf}{facsimile}.
\item Ildefonso Schuster, \textit{The Sacramentary}, trans. Arthur Levelis-Marke, III (London: Burns Oates \& Washbourne, 1927), Sixteenth Sunday, pp.~143--144. Chant reuse is cited as Schuster's comparison.
\item \textit{The Liturgical Year}, continuation under Guéranger's series name, trans. Dom Laurence Shepherd, \textit{Time after Pentecost} II (complete-series XI), second edition (Stanbrook Abbey, Worcester; London: Burns \& Oates, R.~\& T.~Washbourne, and Art \& Book Company, 1909), Sixteenth Sunday, pp.~356--372, especially 359,363,365,370--371; title and preface pp.~iii--iv. \sourceurl{https://archive.org/download/V11TheLiturgicalYear/V11TheLiturgicalYear.pdf}{Facsimile}. The checked preface establishes continuation but does not name its writer.
\item United States Conference of Catholic Bishops, \textit{New American Bible Revised Edition}, official English web text: \sourceurl{https://bible.usccb.org/bible/ephesians/3}{Ephesians~3}, notes on 3:14--21 and 14--15; \sourceurl{https://bible.usccb.org/bible/luke/14}{Luke~14}, note on 14:5; \sourceurl{https://bible.usccb.org/bible/psalms/0}{Psalms introduction}, shared composition boundary. Chapter notes are paraphrased; the introduction's chronology comes through the corpus.
\item Scripture chronology sources: \textit{Catholic Encyclopedia}, ``Epistle to the Ephesians'' (V, 1909), ``St Paul'' (XI, 1911), ``Synoptics'' (XIV, 1912), ``Christmas'' and ``Chronology of the Life of Jesus Christ'' (III, 1908), the corpus's cited composition and Nativity passages; Pontifical Biblical Commission, \textit{De auctore, de tempore compositionis et de historica veritate Evangeliorum secundum Marcum et secundum Lucam}, responses I--IX, AAS~4 (1912); George Leo Haydock, \textit{Douay-Rheims with Haydock Commentary}, Loreto/Feeney Memorial edition (2014), Psalm~70 captivity and Usher-chronology note; Challoner, Ps.~3:1 and Dan.~1:1. These are the distinct chronology assertions projected on page~2.
\item William Shakespeare, \textit{Sonnets}~57, line~5, with lines~1--6,9--14; Folger edition, ed. Barbara A.~Mowat and Paul Werstine, with Michael Poston and Rebecca Niles, FDT~0.9.0.1, created 31 July 2015. \sourceurl{https://www.folger.edu/explore/shakespeares-works/shakespeares-sonnets/read/57/}{Sonnet 57}; \sourceurl{https://folger-main-site-assets.s3.amazonaws.com/uploads/2022/11/shakespeares-sonnets_TXT_FolgerShakespeare.txt}{complete text}, lines~964--979. Shakespeare's underlying English is public domain; Folger's modern editorial material is distinct.
\item Friedrich Nietzsche, \textit{Human, All Too Human: A Book for Free Spirits}, trans. Alexander Harvey (Chicago: Charles H.~Kerr, 1908), ``History of the Moral Feelings,'' aphorism~87, p.~109, with §§82--89; \sourceurl{https://www.gutenberg.org/files/38145/38145-h/38145-h.htm}{Gutenberg eBook 38145}, released 26 November 2011. The offered edition contains its three divisions through aphorism~144.
\item Elizabeth Barrett Browning, \textit{Sonnets from the Portuguese} XLIII, lines~2--3 in the complete sonnet (Caradoc Press, Bedford Park, 1906), David Price transcription, \sourceurl{https://www.gutenberg.org/cache/epub/2002/pg2002.txt}{Gutenberg eBook 2002}, released 1 December 1999, updated 20 August 2026; physical lines~994--1007.
\item Percy Fitzgerald, \textit{Fatal Zero}, in \textit{All the Year Round}, conducted by Charles Dickens, new series I (London, 1869), installment of 16 January, p.~163, left column, paragraph beginning ``Of course he had not heard of my fall in the world.'' \sourceurl{https://upload.wikimedia.org/wikipedia/commons/0/04/All_the_Year_Round_-_Series_2_-_Volume_1.djvu}{Volume facsimile}, digital leaf~173. Authorship: Charles Dickens, letter to Mrs James T.~Fields, Glasgow, 16 December 1868, in \textit{Letters}, ed. Mamie Dickens and Georgina Hogarth, III (London: Chapman and Hall, 1882), \sourceurl{https://www.gutenberg.org/cache/epub/25854/pg25854.txt}{Gutenberg eBook 25854}, released 20 June 2008, physical lines~7797--7798.
\item United States Congress, \textit{Congressional Record---Senate}, 4 October 1977, p.~32287, right column, Riegle's natural-gas speech; \sourceurl{https://www.congress.gov/95/crecb/1977/10/04/GPO-CRECB-1977-pt25-5-2.pdf}{GPO bound-record segment}, PDF p.~90. Doug Batchelor, \textit{How to Keep the Sabbath Holy} (Roseville: Amazing Facts, 2014), ``The Ox in the Ditch,'' p.~78, first paragraph, \sourceurl{https://www.amazingfacts.org/wp-content/uploads/2025/06/BK-HKSH.pdf}{publisher PDF}, p.~77; protected modern corroborant of the idiom only.
\item Anthony of Padua, \textit{Sermones dominicales}, Basilica's official Latin web text, Sunday XVI after Pentecost §§2,7,12 and Sunday XVII §§2,6,8--16: different Pauline/Gospel pairings and the distorted-desire analogue. Triptych, GPT, Tenth, Eleventh, Twelfth, Fourteenth and Fifteenth Sundays after Pentecost, named proposal passages and source-grounded sections cited in the exploratory entries; the bounded precedent audit and exact source paths are in \texttt{research/scope.md}.
\end{itemize}
Online research witnesses were accessed on 5 September 2026; chronology retains its separately recorded corpus source states.
}
